\chapter{上半带的公式，先看它们在节省哪种位置}\label{ch:upper}
\section{这一讲还要证明什么}
上一讲说明，合格的计数表会自动接出剪切和同余输送。本讲只负责证明：对每个允许的 $p,s$，真的有这样的表。我们仍记 $r=p+s$，$1\le s<p$。

不必一开始记住所有公式。先按新值 $D$ 还剩多少份来理解三种选择。新载体保留 $p+s$ 份，因此可用于输出的 $D$ 只有 $p-s$ 份。
\begin{enumerate}
\item $s$ 很大、接近 $p$ 时，$D$ 剩得少，直接设计不使用 $D$ 的输出。
\item $s$ 在中段时，让旧三类取用更均衡，用一个整数 $e$ 在两种输出之间调整。
\item $s$ 很小时，用另外两级相邻参数保证根系数的最大公因数足够小。
\end{enumerate}
三类设计的算术收尾相同，不同的是剩余位置怎么安排。
\begin{proposition}[上半带的充分性]\label{prop:upper}
若 $p\ge11$、$3p<n<4p$，全部合法输入可归零。
\end{proposition}
\begin{proof}
$p\ge83$ 时，由引理\ref{lem:upper-supply}得到定理\ref{thm:terminal}的资源。$11\le p<83$ 且 $n=3p+s$、$s\le9$ 时用引理\ref{lem:small-supply}；其余小参数由引理\ref{lem:upper-small}处理。三者都接命题\ref{prop:entry}的入口及引理\ref{lem:zero-tail}的收尾。
\end{proof}

\section{一个可行整数不够，需要相邻两个}
典型输出写成 $o_e=(\alpha-e,p-3\alpha-e,2\alpha-s+2e,s)$。当 $e$ 增加一，恰好是一份 $A$、一份 $B$ 换成两份 $C$。这就是上一讲产生剪切的相邻输出。

所有库存限制会变成 $L\le e\le U$。我们必须证明 $L+1\le U$，因为一个可行点只给一个反射，两点才给反射之间的比较。

例如 $p=83,s=10$，取 $\alpha=14$。来源 $(28,14,41)$，剩余库存为 $(27,55,11,73)$。可选 $e=-9,-8$，得到输出
\[(23,50,0,10),qquad(22,49,2,10).
\]
两者都非负、和为83、没有超库存。这里 $e$ 为负完全可以：$e$ 是参数，不是直接取用份数；真正不能为负的是输出的四项。

第二个相邻 $\alpha=15$ 同样给两张表。两个 $\alpha$ 互素；根系数 $H=3n\alpha-2p^2$ 分别为 $-2900,-2123$，最大公因数为一。这就满足上一讲的资源条件。
\begin{lemma}[大素数的三组计数公式]\label{lem:upper-supply}
若 $p\ge83$、$p<r<2p$，则存在满足推论\ref{cor:resource}且 $\delta=\gcd(n,2)$ 的相邻表。
\end{lemma}
\begin{proof}
全部取 $q_0=2$。按下列三种范围选计数。

\par\noindent\textbf{第一类：靠近 $4p$。} 若 $r>3p/2$，置 $d=2p-r$，取 $\alpha=\lceil d/2\rceil+2,\lceil d/2\rceil+3$，来源 $(\alpha,0,p-\alpha)$，输出
\[(l,p-\alpha+l,\alpha-2l,0).
\]
其精确区间为 $L\le l\le U$，其中
\[L=\max(0,\lceil(d-\alpha)/2\rceil),\qquad
U=\min(\lfloor\alpha/2\rfloor,p-2\alpha).
\]
取 $l=L+1,L$。有 $L+1\le\lfloor\alpha/2\rfloor$；由 $2\alpha\le d+7$、$L+1\le d/4+1$，另一约束由 $p\ge5d/4+8$ 保证。$d\le(p-1)/2,p\ge83$ 足够。

\par\noindent\textbf{第二类：上半带中部。} 写 $r=p+s$、$10\le s\le(p-1)/2$。取 $\alpha=\lceil p/6\rceil,\lceil p/6\rceil+1$，来源 $(2\alpha,\alpha,p-3\alpha)$，输出
\[(\alpha-e,p-3\alpha-e,2\alpha-s+2e,s).
\]
可行区间为 $L\le e\le U$，
\[L=\max(5\alpha-p,\lceil s/2\rceil-\alpha),\qquad
U=\min(\alpha,2\alpha+s-\lceil p/2\rceil).
\]
取 $e=L,L+1$。$L+1\le U$ 的四个比较是
\begin{align*}
4\alpha&\le p-1,&3\alpha&\le\lfloor p/2\rfloor+s-1,\\
\lceil s/2\rceil+1&\le2\alpha,&\lceil p/2\rceil+1&\le3\alpha+\lfloor s/2\rfloor.
\end{align*}

下面说明四个比较的余量。因为 $p/6\le\alpha\le(p+11)/6$，第一式在 $p\ge25$ 时成立；
第二式左侧至多比 $\lfloor p/2\rfloor$ 大 $6$，而 $s-1\ge9$；
第三式使用 $s\le(p-1)/2$，在 $p\ge15$ 时足够；
第四式由 $3\alpha\ge\lceil p/2\rceil$ 及 $\lfloor s/2\rfloor\ge5$ 成立。
此外 $4\alpha\le p$ 保证第一类来源不超量，$6\alpha\ge p$ 保证载体来源容量，
$s\le p/2$ 保证新值 $D$ 的库存。因此所列整数区间确实包含相邻两个输出。

\par\noindent\textbf{第三类：紧邻 $3p$ 的九个偏移。} 若 $1\le s\le9$，令 $a_0=\lfloor s/3\rfloor+1$，取 $\alpha=a_0,a_0+1$、$h=\lfloor p/(3\alpha)\rfloor-1,\lfloor p/(3\alpha)\rfloor$。令 $t=p-3\alpha h$，来源和输出为
\[
(\alpha(h+1),\alpha h,p-\alpha(2h+1)),\quad
(\alpha h-e,p-\alpha(2h+1)-e,(\alpha-s)h+\alpha+2e,sh).
\]
输出区间为
\[
L=\max(0,2\alpha-t,\lceil((s-\alpha)h-\alpha)/2\rceil),\quad
U=\min(\alpha h,\alpha h+t-\alpha,\lfloor(sh+\alpha+s-t)/2\rfloor).
\]
可取 $e=L,L+1$，事实上 $L+2\le U$。逐项移项，充分条件是
\[
\begin{gathered}
h\ge4,\quad(3\alpha-s)h\ge\max(s,\alpha+4),\\
sh\ge\max(5\alpha+3-s,3\alpha+4-s),\\
\alpha h\ge\max(4\alpha+4-s,3\alpha+2).
\end{gathered}
\]
对 $\alpha=1,2,3,4,5$，所需 $h$ 下界依次为 $7,12,8,4,4$；$p\ge83$ 时实际最小值为 $26,12,8,5,4$。这同时保证来源及 $sh\le p-s$ 的库存。


\par\noindent\textbf{三类计数的最大公因子。}
到这里已经有真实反射，但还需把它们合成足够浅的根与同余群。
这一步只计算前一讲的 $H=n(2j+\alpha)-2p^2$。
三种公式均满足迹零条件。第一种的 $H=n\alpha-2p^2$，第二种为 $3n\alpha-2p^2$，相邻 $\alpha$ 给 $\gcd(H)=\gcd(n,2)$。第三种相邻 $h$ 的差为 $2\alpha n$，相邻 $\alpha$ 互素；再用 $H\equiv-2p^2\pmod n$ 及奇偶性得到同一结论。所有横向系数含两个相邻整数。
\end{proof}

\section{83是一个统一够用的界，不是数学相变点}
三套计数都可以用 $p\ge83$ 一次保证，但其中两套在小素数上早已能用。偏移 $s\ge10$ 的情形，只需回到四个容量不等式逐项检查，小参数就能补齐。
\begin{lemma}[偏移至少10时的小素数]\label{lem:upper-small}
若 $p\ge11$、$10\le s<p$，引理\ref{lem:upper-supply}的前两种计数已提供推论\ref{cor:resource}的全部资源。
\end{lemma}
\begin{proof}
中部 $s\le(p-1)/2$ 时，四个容量比较对 $p\ge29$ 成立：第一项只需 $p\ge25$；第二项由 $s\ge10$ 保证；第三项只需 $p\ge15$；第四项由 $3\alpha\ge\lceil p/2\rceil$ 保证。唯一较小素数是 $p=23,s=10,11$，此时 $\alpha=4,5$，直接代入四式均成立。
靠近 $4p$ 时，$d=p-s\le(p-1)/2$。充分条件 $p\ge5d/4+8$ 对 $p\ge23$ 成立。若 $p=11,13,17$，还可用 $d\le p-10$，分别得到右侧上界 $9.25,11.75,16.75$。若 $p=19$，$d\le8$ 同样成立；仅余 $d=9$，其 $\alpha=7,8$ 均有 $L=1,U=3$。故两张相邻输出仍在库存内。根和横向系数的最大公因子计算与引理\ref{lem:upper-supply}相同。
\end{proof}
这个引理的作用是缩小真正要查表的区域。不能因为某个统一证明选了83，就把所有小于83的参数都当作没有结构的特殊例子。

\section{直接造出n倍剪切：用一次守恒关系}
记来源为 $(i,j,k)$，新均值满足 $iA+jB+kC=pD$，全局还满足 $pA+pB+rC=0$。如果一个输出组的四项计数允许增加
\[\Delta_t=(ti-p, tj-p, tk-r, n-tp),
\]
那么项数不变，而输出总和增加
\[
t(iA+jB+kC)-(pA+pB+rC)+(n-tp)D=nD.
\]
所以新第一块均值增加 $nD/p=-nz'$，载体 $z'$ 不变。两条返回比较，恰好是 $(a',z')\mapsto(a'-nz',z')$，即 $U(-n)$。

这个根参数 $n$ 直接来自总和守恒。唯一还需检查的是：移动前后的两张输出表都不能超库存。五返回公式就是把这项检查与两个横向参数一起安排。
\begin{lemma}[九个低偏移的有限补充]\label{lem:small-supply}
若 $11\le p<83$、$1\le s\le9$，则全部合法 $K_{p,p+s}$ 可归零。
\end{lemma}
\begin{proof}
先给五返回公式。令 $a=\lceil s/2\rceil$，对 $\alpha=a,a+1$ 写 $p=3\alpha d+\rho$，$0\le\rho<3\alpha$。来源为 $(\alpha(d+1),\alpha d,\alpha d+\rho-\alpha)$，输出为
\[
o_e=(\alpha d-e,\alpha d+\rho-\alpha-e,(\alpha-s)d+\alpha+2e,sd).
\]
对于第一组 $\alpha=a$，还使用 $o_e+\Delta$，其中 $\Delta=(3\alpha-\rho,-\rho,2\rho-3\alpha-s,s)$。可同时容纳 $o_e,o_{e+1},o_e+\Delta$ 的区间是
\begin{align*}
L={}&\max\left(0,5\alpha-2\rho,
\left\lceil\frac{(s-\alpha)d-\alpha}{2}\right\rceil,
\left\lceil\frac{(s-\alpha)d+2\alpha+s-2\rho}{2}\right\rceil\right),\\
U={}&\min\left(\alpha d-\alpha-\mathbf1_{\rho=0},
\left\lfloor\frac{sd+\alpha+s-\rho-2}{2}\right\rfloor,
\left\lfloor\frac{sd+4\alpha+2s-3\rho}{2}\right\rfloor\right).
\end{align*}
逐项比较给充分条件 $d\ge7,sd\ge9\alpha+2-s,(3\alpha-s)d\ge4\alpha+s+2$。例如最紧的三式是 $5\alpha\le\alpha d-\alpha-1$、$10\alpha\le sd+\alpha+s-2$、$(s-\alpha)d+2\alpha+s\le2(\alpha d-\alpha-1)$；其余比较的余量至少分别为 $\alpha d+2\alpha+s-\rho-3$ 和 $\alpha d+2\alpha+s-\rho-1$。第二组只需相邻输出，条件降为 $d\ge4,sd\ge3\alpha+2-s$。由此得到九个充分截点
\[
T_s=(42,24,42,42,63,63,84,84,105).
\]
当 $p\ge T_s$ 时，以上计数直接可用。守恒关系 $pA+pB+rC=0$ 使沿 $\Delta$ 的输出和增加 $nD$，故两返回之商为 $U(-n)$。相邻倍率扩展为 $U(nR)$，两横向系数 $a,a+1$ 由引理\ref{lem:transverse}给 $\Gamma(n^2,R)$，接定理\ref{thm:terminal}。


\par\noindent\textbf{有限剩余。}
余下的114组参数及233对相邻表完整列在附录\ref{app:certificates}。它们的非负库存、迹零等式、$\alpha$ 的最大公因子和 $H$ 的最大公因子由命题\ref{prop:root-certificate}核对。除 $(p,s)=(11,3)$ 外均满足推论\ref{cor:resource}。

\par\noindent\textbf{唯一单列参数 $(p,s)=(11,3)$。}
此时 $n=36$，一般最大公因子接口改由两组明确矩阵完成。
在此例，来源 $(4,2,5)$、输出 $(2,5,1,3)$ 与 $(3,0,2,6)$ 给
\[
J=\begin{pmatrix}-27&568\\-2&27\end{pmatrix},\quad
M=\begin{pmatrix}45&-404\\-2&27\end{pmatrix},\quad MJ^{-1}=U(-36).
\]
库存为 $(3,7,4,8)$，两输出均可用；$J^2=-407I$。相邻输出 $(1,4,3,3)$ 的反射 $G$ 满足
\[J^{-1}G=\begin{pmatrix}1&-972/37\\0&-35/37\end{pmatrix}.
\]
故 $R_0=\Z[1/(5\cdot7\cdot37)]$ 上有 $U(36R_0)$。另一个表 $(4,1,6;2,7,1,1)$ 给 $K=\left(\begin{smallmatrix}-52&905\\-3&52\end{smallmatrix}\right)$、$K^2=-11I$。横向系数 $2,3$ 合并给 $\Gamma(36^2,R_0)$，再取其子群 $\Gamma(4\cdot36^2,R_0)$ 接定理\ref{thm:terminal}。
\end{proof}

\section{剩余114组表，承担的是哪一个有限命题}
到此，无穷多个大 $p$ 已由公式覆盖，小 $p$ 的不少参数也由五返回公式覆盖。还保留114组物理根资源。它们各列一对或几对具体计数，检查非负、总数、迹零以及两个最大公因数。

查表以后，对任意有理参数的终止仍由同余输送证明。表里不枚举数值高度，不说“随机一千次都成功所以全部成功”。附录保留所有计数，你可以用上一讲三十四元的方式逐行读它。
